% Diagrama da Metodologia: caminhos de uma leitura na hierarquia do gfx803 e
% o que cada experimento mede. Incluído em main.tex por % Diagrama da Metodologia: caminhos de uma leitura na hierarquia do gfx803 e
% o que cada experimento mede. Incluído em main.tex por % Diagrama da Metodologia: caminhos de uma leitura na hierarquia do gfx803 e
% o que cada experimento mede. Incluído em main.tex por % Diagrama da Metodologia: caminhos de uma leitura na hierarquia do gfx803 e
% o que cada experimento mede. Incluído em main.tex por \input{fig_hierarquia}.
\begin{tikzpicture}[
  font=\footnotesize,
  nivel/.style={draw, rounded corners=2pt, minimum height=0.9cm,
                minimum width=2.3cm, align=center, fill=black!5},
  seta/.style={-latex, semithick},
  node distance=1.25cm,
]
\node[nivel] (wi) {\textit{Work-item}\\registradores};
\node[nivel, right=of wi] (l1) {L1 vetorial\\16~KB por CU};
\node[nivel, right=of l1] (l2) {L2\\2~MB};
\node[nivel, right=of l2] (vr) {VRAM\\GDDR5, 8~GB};
\draw[seta] (wi) -- node[above=1pt, font=\scriptsize]{\texttt{GLC=0}} (l1);
\draw[seta] (l1) -- (l2);
\draw[seta] (l2) -- (vr);
\draw[seta] (wi.north) to[out=20, in=160]
  node[above=-1pt]{\texttt{glc}: contorna o L1 (Exp.~1, 3 e 4)} (l2.north);
\draw[seta, dashed] (wi.south) to[out=-14, in=-166]
  node[below=-1pt]{\texttt{slc}: acesso direto à VRAM? (Exp.~2)}
  (vr.south);
\end{tikzpicture}
.
\begin{tikzpicture}[
  font=\footnotesize,
  nivel/.style={draw, rounded corners=2pt, minimum height=0.9cm,
                minimum width=2.3cm, align=center, fill=black!5},
  seta/.style={-latex, semithick},
  node distance=1.25cm,
]
\node[nivel] (wi) {\textit{Work-item}\\registradores};
\node[nivel, right=of wi] (l1) {L1 vetorial\\16~KB por CU};
\node[nivel, right=of l1] (l2) {L2\\2~MB};
\node[nivel, right=of l2] (vr) {VRAM\\GDDR5, 8~GB};
\draw[seta] (wi) -- node[above=1pt, font=\scriptsize]{\texttt{GLC=0}} (l1);
\draw[seta] (l1) -- (l2);
\draw[seta] (l2) -- (vr);
\draw[seta] (wi.north) to[out=20, in=160]
  node[above=-1pt]{\texttt{glc}: contorna o L1 (Exp.~1, 3 e 4)} (l2.north);
\draw[seta, dashed] (wi.south) to[out=-14, in=-166]
  node[below=-1pt]{\texttt{slc}: acesso direto à VRAM? (Exp.~2)}
  (vr.south);
\end{tikzpicture}
.
\begin{tikzpicture}[
  font=\footnotesize,
  nivel/.style={draw, rounded corners=2pt, minimum height=0.9cm,
                minimum width=2.3cm, align=center, fill=black!5},
  seta/.style={-latex, semithick},
  node distance=1.25cm,
]
\node[nivel] (wi) {\textit{Work-item}\\registradores};
\node[nivel, right=of wi] (l1) {L1 vetorial\\16~KB por CU};
\node[nivel, right=of l1] (l2) {L2\\2~MB};
\node[nivel, right=of l2] (vr) {VRAM\\GDDR5, 8~GB};
\draw[seta] (wi) -- node[above=1pt, font=\scriptsize]{\texttt{GLC=0}} (l1);
\draw[seta] (l1) -- (l2);
\draw[seta] (l2) -- (vr);
\draw[seta] (wi.north) to[out=20, in=160]
  node[above=-1pt]{\texttt{glc}: contorna o L1 (Exp.~1, 3 e 4)} (l2.north);
\draw[seta, dashed] (wi.south) to[out=-14, in=-166]
  node[below=-1pt]{\texttt{slc}: acesso direto à VRAM? (Exp.~2)}
  (vr.south);
\end{tikzpicture}
.
\begin{tikzpicture}[
  font=\footnotesize,
  nivel/.style={draw, rounded corners=2pt, minimum height=0.9cm,
                minimum width=2.3cm, align=center, fill=black!5},
  seta/.style={-latex, semithick},
  node distance=1.25cm,
]
\node[nivel] (wi) {\textit{Work-item}\\registradores};
\node[nivel, right=of wi] (l1) {L1 vetorial\\16~KB por CU};
\node[nivel, right=of l1] (l2) {L2\\2~MB};
\node[nivel, right=of l2] (vr) {VRAM\\GDDR5, 8~GB};
\draw[seta] (wi) -- node[above=1pt, font=\scriptsize]{\texttt{GLC=0}} (l1);
\draw[seta] (l1) -- (l2);
\draw[seta] (l2) -- (vr);
\draw[seta] (wi.north) to[out=20, in=160]
  node[above=-1pt]{\texttt{glc}: contorna o L1 (Exp.~1, 3 e 4)} (l2.north);
\draw[seta, dashed] (wi.south) to[out=-14, in=-166]
  node[below=-1pt]{\texttt{slc}: acesso direto à VRAM? (Exp.~2)}
  (vr.south);
\end{tikzpicture}
